\chapter*{Решения}
\addcontentsline{toc}{chapter}{Решения}
\markboth{Решения}{Решения}

\section*{Тема 1}

\begin{proofbox}[Решение на Задача 1.1]
Твърдим, че $\mathcal A = \big\{\bigcup_{i\in S} B_i : S\subseteq
\{1,\dots,k\}\big\}$ (обединение по всяко подмножество на индексите,
като за $S=\varnothing$ приемаме обединението за $\varnothing$).
Тази колекция е $\sigma$-алгебра: съдържа $\Omega$ (за $S=\{1,\dots,
k\}$); затворена е относно допълнение, защото допълнението на
$\bigcup_{i\in S}B_i$ е точно $\bigcup_{i\notin S}B_i$ (понеже
блоковете покриват $\Omega$ и~са несъвместими); и~е затворена относно
(крайно) обединение, защото обединение на такива множества отговаря
на обединението на съответните индексни множества. Освен това всяка
$\sigma$-алгебра, съдържаща $B_1,\dots,B_k$, трябва да съдържа всички
такива обединения по затвореност относно обединение, така че
$\mathcal A$ е точно породената $\sigma$-алгебра. Броят на
подмножествата на $\{1,\dots,k\}$ е $2^k$, а~различните подмножества
дават различни (несъвпадащи) обединения, понеже блоковете са
непразни и~по двойки несъвместими. Следователно $|\mathcal A|=2^k$.
\end{proofbox}

\begin{proofbox}[Решение на Задача 1.2]
Проверяваме трите аксиоми за $\mathcal A=\bigcap_{i\in I}\mathcal
A_i$. \textbf{1)} $\Omega\in\mathcal A_i$ за всяко $i$ (всяко
$\mathcal A_i$ е $\sigma$-алгебра), следователно $\Omega\in\mathcal
A$. \textbf{2)} Нека $A\in\mathcal A$, т.е.\ $A\in\mathcal A_i$ за
всяко $i\in I$. За всяко $i$, понеже $\mathcal A_i$ е затворена
относно допълнение, $A^c\in\mathcal A_i$. Това важи за всяко $i$,
следователно $A^c\in\bigcap_i\mathcal A_i=\mathcal A$. \textbf{3)}
Нека $A_1,A_2,\dots\in\mathcal A$, т.е.\ за всяко $n$ и~всяко $i$:
$A_n\in\mathcal A_i$. Фиксираме $i\in I$: понеже всички $A_n\in
\mathcal A_i$ и~$\mathcal A_i$ е затворена относно изброимо
обединение, $\bigcup_n A_n\in\mathcal A_i$. Това важи за всяко
$i\in I$, следователно $\bigcup_n A_n\in\bigcap_i\mathcal A_i=
\mathcal A$. \qedbg
\end{proofbox}

\begin{proofbox}[Решение на Задача 1.3]
Понеже $\mathcal A$ е затворена относно допълнение, $B^c\in\mathcal
A$. По закона на Де Морган, $A\cap B^c = (A^c\cup B)^c$. Понеже
$A^c,B\in\mathcal A$, тяхното обединение $A^c\cup B\in\mathcal A$
(изброимо обединение с~всички останали членове $\varnothing$), а~оттам и~допълнението му $(A^c\cup B)^c=A\cap B^c\in\mathcal A$ (пак
по затвореност относно допълнение). Но $A\setminus B = A\cap B^c$,
следователно $A\setminus B\in\mathcal A$. Аналогично, $B\setminus A
= B\cap A^c\in\mathcal A$. Накрая, $A\triangle B = (A\setminus B)
\cup(B\setminus A)$ е обединение на две множества от $\mathcal A$,
следователно $A\triangle B\in\mathcal A$. \qedbg
\end{proofbox}

\section*{Тема 2}

\begin{proofbox}[Решение на Задача 2.1]
От известната формула $\mathbb P(A\cup B)=\mathbb P(A)+\mathbb P(B)-
\mathbb P(A\cap B)$ и~$\mathbb P(A\cup B)\le1$ (защото $A\cup B
\subseteq\Omega$ и~монотонност), получаваме
\[
  \mathbb P(A)+\mathbb P(B)-\mathbb P(A\cap B) \le 1 \quad
  \Longrightarrow\quad \mathbb P(A\cap B) \ge \mathbb P(A)+\mathbb
  P(B)-1. \qedbg
\]
\end{proofbox}

\begin{proofbox}[Решение на Задача 2.2]
База ($n=2$): вече доказано в~Тема~2. Индукционна стъпка: нека
твърдението е вярно за $n$ събития и~нека $A_1,\dots,A_{n+1}$ са
произволни. Прилагайки случая за две събития към $\bigcup_{i=1}^n
A_i$ и~$A_{n+1}$:
\[
  \mathbb P\Big(\bigcup_{i=1}^{n+1}A_i\Big) = \mathbb P\Big(\Big(
  \bigcup_{i=1}^n A_i\Big)\cup A_{n+1}\Big) \le \mathbb P\Big(
  \bigcup_{i=1}^n A_i\Big) + \mathbb P(A_{n+1}) \overset{\text{И.П.}}
  \le \sum_{i=1}^n \mathbb P(A_i) + \mathbb P(A_{n+1}) = \sum_{i=1}^{n+1}
  \mathbb P(A_i),
\]
където използвахме индукционното предположение (И.П.) в~предпоследната
стъпка. \qedbg
\end{proofbox}

\begin{proofbox}[Решение на Задача 2.3]
Понеже $\mathbb P(A_n)=1$, имаме $\mathbb P(A_n^c)=1-\mathbb P(A_n)=0$
за всяко $n$. Нека $C_n=\bigcap_{k=1}^n A_k$ (намаляваща редица).
$C_n^c=\bigcup_{k=1}^n A_k^c$, и~по Задача~2.2 (крайна субадитивност):
$\mathbb P(C_n^c)\le\sum_{k=1}^n\mathbb P(A_k^c)=0$, т.е.\ $\mathbb
P(C_n)=1$ за всяко $n$. Понеже $\bigcap_{n=1}^{\infty}A_n=
\bigcap_{n=1}^{\infty}C_n$ и~$(C_n)$ намалява, по непрекъснатостта
на вероятността отгоре (Тема~2):
\[
  \mathbb P\Big(\bigcap_{n=1}^{\infty}A_n\Big) = \lim_{n\to\infty}
  \mathbb P(C_n) = 1. \qedbg
\]
\end{proofbox}

\section*{Тема 3}

\begin{proofbox}[Решение на Задача 3.1]
Условното пространство се състои от $3\times6=18$ равновъзможни
изхода (първи зар $\in\{2,4,6\}$, втори зар произволен). За всяка
стойност на първия зар има точно едно число за втория зар, което
дава сбор~$7$: $(2,5),(4,3),(6,1)$. Следователно
\[
  \mathbb P(\text{сбор}=7 \mid \text{първи зар четен}) = \frac3{18}
  = \frac16.
\]
\end{proofbox}

\begin{proofbox}[Решение на Задача 3.2]
\textbf{1)} $Q(B)=\mathbb P(B\cap A)/\mathbb P(A)\ge0$, защото
числителят и~знаменателят са неотрицателни. \textbf{2)} $Q(\Omega)=
\mathbb P(\Omega\cap A)/\mathbb P(A) = \mathbb P(A)/\mathbb P(A)=1$.
\textbf{3)} Нека $B_1,B_2,\dots$ са по двойки несъвместими. Тогава
$(B_i\cap A)$ също са по двойки несъвместими (наследяват
несъвместимостта на $B_i$), следователно по $\sigma$-адитивност на
$\mathbb P$:
\[
  Q\Big(\bigcup_i B_i\Big) = \frac{\mathbb P\big((\bigcup_iB_i)\cap A
  \big)}{\mathbb P(A)} = \frac{\mathbb P\big(\bigcup_i(B_i\cap A)
  \big)}{\mathbb P(A)} = \frac{\sum_i \mathbb P(B_i\cap A)}{\mathbb
  P(A)} = \sum_i \frac{\mathbb P(B_i\cap A)}{\mathbb P(A)} = \sum_i
  Q(B_i).
\]
И трите аксиоми са изпълнени, следователно $Q$ е вероятностна мярка.
\qedbg
\end{proofbox}

\begin{proofbox}[Решение на Задача 3.3]
Нека $M_i$ = „детайлът е от машина $i$“, $D$ = „детайлът е бракуван“.
$\mathbb P(M_1)=0.3$, $\mathbb P(M_2)=0.45$, $\mathbb P(M_3)=0.25$;
$\mathbb P(D\mid M_1)=0.02$, $\mathbb P(D\mid M_2)=0.03$, $\mathbb
P(D\mid M_3)=0.05$. По формулата за пълната вероятност:
\[
  \mathbb P(D) = 0.3\cdot0.02+0.45\cdot0.03+0.25\cdot0.05 = 0.006+
  0.0135+0.0125 = 0.032.
\]
По формулата на Бейс:
\[
  \mathbb P(M_2\mid D) = \frac{\mathbb P(D\mid M_2)\,\mathbb P(M_2)}
  {\mathbb P(D)} = \frac{0.0135}{0.032} = \frac{27}{64} \approx
  0.422.
\]
\end{proofbox}

\section*{Тема 4}

\begin{proofbox}[Решение на Задача 4.1]
За произволно $x\in\mathbb R$: $\max(X,Y)<x \iff X<x \text{ и~} Y<x$,
т.е.\ $\{Z<x\}=\{X<x\}\cap\{Y<x\}$. И двете множества са в~$\mathcal
A$ (защото $X,Y$ са случайни величини), следователно сечението им
$\{Z<x\}\in\mathcal A$ (затвореност относно допълнение и~обединение,
приложена за да получим затвореност относно сечение). Понеже това е
вярно за всяко $x$, $Z=\max(X,Y)$ е случайна величина. \qedbg
\end{proofbox}

\begin{proofbox}[Решение на Задача 4.2]
От свойствата на индикаторната функция (Следствие 4.1), $I_A\cdot I_B
= I_{A\cap B}$. Следователно, използвайки $\mathbb EI_C=\mathbb P(C)$
за произволно събитие $C$:
\[
  \mathbb EI_A + \mathbb EI_B - \mathbb E(I_AI_B) = \mathbb P(A) +
  \mathbb P(B) - \mathbb E I_{A\cap B} = \mathbb P(A)+\mathbb P(B)-
  \mathbb P(A\cap B) = \mathbb P(A\cup B),
\]
където последното равенство е добре известната формула за
обединение на две събития (Тема~2). \qedbg
\end{proofbox}

\begin{proofbox}[Решение на Задача 4.3]
Нека $u,v\in\mathbb R$ са произволни (в~образите на $g(X)$, $h(Y)$
съответно), и~нека $A=\{x : g(x)=u\}$, $B=\{y:h(y)=v\}$ (в~носителите на $X,Y$). Тогава $\{g(X)=u\}=\bigcup_{x_i\in A}\{X=x_i\}$
и~$\{h(Y)=v\}=\bigcup_{y_j\in B}\{Y=y_j\}$. Използвайки независимостта
на $X,Y$:
\begin{align*}
  \mathbb P\big(g(X)=u,\,h(Y)=v\big) &= \sum_{x_i\in A}\sum_{y_j\in
  B} \mathbb P(X=x_i,Y=y_j) = \sum_{x_i\in A}\sum_{y_j\in B} \mathbb
  P(X=x_i)\,\mathbb P(Y=y_j)\\
  &= \Big(\sum_{x_i\in A}\mathbb P(X=x_i)\Big)\Big(\sum_{y_j\in B}
  \mathbb P(Y=y_j)\Big) = \mathbb P(X\in A)\,\mathbb P(Y\in B)\\
  &= \mathbb P\big(g(X)=u\big)\,\mathbb P\big(h(Y)=v\big).
\end{align*}
Понеже $u,v$ бяха произволни, $g(X)\perp\!\!\!\perp h(Y)$. \qedbg
\end{proofbox}

\section*{Тема 5}

\begin{proofbox}[Решение на Задача 5.1]
$DX=\mathbb E(X-\mathbb EX)^2$, а~$(X-\mathbb EX)^2\ge0$ във всяка
точка, следователно по монотонност на очакването $\mathbb E(X-
\mathbb EX)^2\ge \mathbb E0=0$, т.е.\ $DX\ge0$. За второто:
\[
  D(aX+b) = \mathbb E\big(aX+b-\mathbb E(aX+b)\big)^2 = \mathbb E
  \big(aX+b-a\,\mathbb EX-b\big)^2 = \mathbb E\big(a(X-\mathbb EX)
  \big)^2 = a^2\,\mathbb E(X-\mathbb EX)^2 = a^2DX. \qedbg
\]
\end{proofbox}

\begin{proofbox}[Решение на Задача 5.2]
$g_X(s)=(ps+1-p)^n$. Първа производна: $g_X'(s)=np(ps+1-p)^{n-1}$,
при $s=1$: $g_X'(1)=np(p+1-p)^{n-1}=np$. Втора производна:
$g_X''(s)=n(n-1)p^2(ps+1-p)^{n-2}$, при $s=1$: $g_X''(1)=n(n-1)p^2$.
Тогава
\[
  DX = g_X''(1)+g_X'(1)-(g_X'(1))^2 = n(n-1)p^2+np-n^2p^2 = np-np^2
  = np(1-p),
\]
което съвпада точно с~познатите $\mathbb EX=np$, $DX=np(1-p)$.
\end{proofbox}

\begin{proofbox}[Решение на Задача 5.3]
\[
  g_X(s) = \sum_{k=0}^n s^k\,\mathbb P(X=k) = \sum_{k=0}^n s^k\binom
  nk p^k(1-p)^{n-k} = \sum_{k=0}^n \binom nk (ps)^k(1-p)^{n-k}
  \overset{\text{бином на Нютон}}{=} (ps+1-p)^n. \qedbg
\]
\end{proofbox}

\section*{Тема 6}

\begin{proofbox}[Решение на Задача 6.1]
За $X\sim Geom(p)$: $\mathbb P(X>k)=(1-p)^k$ за $k\ge0$ (стандартна
опашкова формула). Понеже $m+n\ge m$, $\{X>m+n\}\subseteq\{X>m\}$,
следователно
\[
  \mathbb P(X>m+n\mid X>m) = \frac{\mathbb P(X>m+n,\,X>m)}{\mathbb
  P(X>m)} = \frac{\mathbb P(X>m+n)}{\mathbb P(X>m)} = \frac{(1-p)^{m+n}}
  {(1-p)^m} = (1-p)^n = \mathbb P(X>n). \qedbg
\]
\end{proofbox}

\begin{proofbox}[Решение на Задача 6.2]
За $0\le k\le n+m$, по формулата за пълна вероятност (условие по
стойността на $X$) и~независимостта на $X,Y$:
\begin{align*}
  \mathbb P(X+Y=k) &= \sum_{i=0}^{k} \mathbb P(X=i)\,\mathbb P(Y=k-i)
  = \sum_{i=0}^k \binom ni p^i(1-p)^{n-i}\binom m{k-i}p^{k-i}(1-p)^{
  m-k+i}\\
  &= p^k(1-p)^{n+m-k}\sum_{i=0}^k\binom ni\binom m{k-i}
  \overset{\text{Вандермонд}}{=} p^k(1-p)^{n+m-k}\binom{n+m}k,
\end{align*}
т.е.\ точно вероятностната функция на $Bin(n+m,p)$. \qedbg
\end{proofbox}

\begin{proofbox}[Решение на Задача 6.3]
$\lambda=np=1000\cdot0.002=2$. По теоремата на Поасон, $X\approx
Poisson(2)$, следователно
\[
  \mathbb P(X=3) \approx e^{-2}\frac{2^3}{3!} = \frac{4e^{-2}}3
  \approx 0.180.
\]
Тъй като $n$ е голямо, а~$p$ - малко (с~умерено $np$), приближението
е добро.
\end{proofbox}

\section*{Тема 7}

\begin{proofbox}[Решение на Задача 7.1]
\begin{align*}
  \operatorname{cov}(aX+b,cY+d) &= \mathbb E\big[(aX+b)(cY+d)\big] -
  \mathbb E(aX+b)\,\mathbb E(cY+d)\\
  &= ac\,\mathbb EXY+ad\,\mathbb EX+bc\,\mathbb EY+bd -
  (a\,\mathbb EX+b)(c\,\mathbb EY+d)\\
  &= ac\,\mathbb EXY+ad\,\mathbb EX+bc\,\mathbb EY+bd - ac\,\mathbb
  EX\,\mathbb EY-ad\,\mathbb EX-bc\,\mathbb EY-bd\\
  &= ac\big(\mathbb EXY-\mathbb EX\,\mathbb EY\big) = ac\cdot
  \operatorname{cov}(X,Y). \qedbg
\end{align*}
\end{proofbox}

\begin{proofbox}[Решение на Задача 7.2]
\begin{align*}
  D(X+Y) &= \mathbb E\big(X+Y-\mathbb E(X+Y)\big)^2 = \mathbb E
  \big((X-\mathbb EX)+(Y-\mathbb EY)\big)^2\\
  &= \mathbb E(X-\mathbb EX)^2 + 2\,\mathbb E\big[(X-\mathbb EX)(Y-
  \mathbb EY)\big] + \mathbb E(Y-\mathbb EY)^2\\
  &= DX + 2\operatorname{cov}(X,Y) + DY. \qedbg
\end{align*}
\end{proofbox}

\begin{proofbox}[Решение на Задача 7.3]
$\mathbb EX = \tfrac13(-1+0+1)=0$. $\mathbb EY=\mathbb EX^2=\tfrac13
(1+0+1)=\tfrac23$. $\mathbb EXY=\mathbb EX^3=\tfrac13(-1+0+1)=0$.
Следователно $\operatorname{cov}(X,Y)=\mathbb EXY-\mathbb EX\,
\mathbb EY = 0-0\cdot\tfrac23=0$. Но $X,Y$ не са независими: $\mathbb
P(X=1,Y=0)=0$ (защото $X=1\Rightarrow Y=X^2=1\ne0$), докато $\mathbb
P(X=1)\,\mathbb P(Y=0)=\tfrac13\cdot\tfrac13=\tfrac19\ne0$.
Следователно нулева ковариация не влече независимост. \qedbg
\end{proofbox}

\section*{Тема 8}

\begin{proofbox}[Решение на Задача 8.1]
Означаваме $A_j=\{Y=y_j\}$, така че $\mathbb E[X\mid Y]=\sum_j
\frac{\mathbb EXI_{A_j}}{\mathbb P(A_j)}I_{A_j}$. Тогава
\[
  \mathbb E\big[X\cdot\mathbb E[X\mid Y]\big] = \sum_j\frac{\mathbb
  EXI_{A_j}}{\mathbb P(A_j)}\,\mathbb E[XI_{A_j}] = \sum_j\frac{
  (\mathbb EXI_{A_j})^2}{\mathbb P(A_j)}.
\]
От друга страна, понеже $I_{A_j}I_{A_k}=0$ за $j\ne k$ (несъвместими)
и~$I_{A_j}^2=I_{A_j}$:
\[
  \mathbb E\big[(\mathbb E[X\mid Y])^2\big] = \mathbb E\Big[\sum_j
  \Big(\frac{\mathbb EXI_{A_j}}{\mathbb P(A_j)}\Big)^2 I_{A_j}\Big] =
  \sum_j \Big(\frac{\mathbb EXI_{A_j}}{\mathbb P(A_j)}\Big)^2 \mathbb
  P(A_j) = \sum_j\frac{(\mathbb EXI_{A_j})^2}{\mathbb P(A_j)}.
\]
Двата израза съвпадат, следователно $\mathbb E[X\cdot\mathbb E[X\mid
Y]] = \mathbb E[(\mathbb E[X\mid Y])^2]$. \qedbg
\end{proofbox}

\begin{proofbox}[Решение на Задача 8.2]
При фиксирано $N=n$, $X\mid N{=}n \sim Bin(n,\tfrac12)$, следователно
$\mathbb E[X\mid N]=\tfrac N2$. По кула-свойството:
\[
  \mathbb EX = \mathbb E\big[\mathbb E[X\mid N]\big] = \mathbb E
  \Big[\frac N2\Big] = \frac12\,\mathbb EN = \frac12\cdot\frac{1+2+
  3+4+5+6}6 = \frac12\cdot\frac{21}6 = \frac74.
\]
\end{proofbox}

\begin{proofbox}[Решение на Задача 8.3]
$\mathbb P(Y=0)=0.1+0.3=0.4$, $\mathbb P(Y=1)=0.2+0.4=0.6$.
\[
  \mathbb E[X\mid Y=0] = 0\cdot\mathbb P(X=0\mid Y=0) + 1\cdot
  \mathbb P(X=1\mid Y=0) = \frac{\mathbb P(X=1,Y=0)}{\mathbb P(Y=0)}
  = \frac{0.3}{0.4} = 0.75,
\]
\[
  \mathbb E[X\mid Y=1] = \frac{\mathbb P(X=1,Y=1)}{\mathbb P(Y=1)} =
  \frac{0.4}{0.6} = \frac23.
\]
\end{proofbox}

\section*{Тема 9}

\begin{proofbox}[Решение на Задача 9.1]
Нормировка: $\int_{-1}^1 c(1-x^2)\,dx = c\big[x-\tfrac{x^3}3\big]_{-1}
^1 = c\cdot\tfrac43=1 \Rightarrow c=\tfrac34$. За $x\in(-1,1)$:
\[
  F_X(x) = \int_{-1}^x \frac34(1-t^2)\,dt = \frac34\Big(x-\frac{x^3}
  3+\frac23\Big) = \frac34x-\frac14x^3+\frac12,
\]
($F_X(x)=0$ за $x\le-1$, $F_X(x)=1$ за $x\ge1$). Тогава $F_X(0)=
\tfrac12$, $F_X(\tfrac12)=\tfrac38-\tfrac1{32}+\tfrac12=\tfrac{27}
{32}$, следователно
\[
  \mathbb P(0<X<\tfrac12) = F_X(\tfrac12)-F_X(0) = \frac{27}{32}-
  \frac{16}{32} = \frac{11}{32}.
\]
\end{proofbox}

\begin{proofbox}[Решение на Задача 9.2]
$g(x)=-\ln x$ е строго намаляваща и~диференцируема върху $(0,1)$, с~обратна $g^{-1}(y)=e^{-y}$ за $y>0$, $(g^{-1}(y))'=-e^{-y}$. По
теоремата за смяна на променливите:
\[
  f_Y(y) = f_X\big(e^{-y}\big)\cdot\big|-e^{-y}\big| = 1\cdot e^{-y}
  = e^{-y}, \qquad y>0
\]
(използвахме $f_X\equiv1$ на $(0,1)$, а~$e^{-y}\in(0,1)$ за $y>0$).
Това е точно плътността на $Exp(1)$, следователно $Y\sim Exp(1)$.
\qedbg
\end{proofbox}

\begin{proofbox}[Решение на Задача 9.3]
\[
  \int_0^{\infty} \mathbb P(X>x)\,dx = \int_0^{\infty}\Big(\int_x^{
  \infty} f_X(t)\,dt\Big)dx.
\]
Областта на интегриране е $\{(x,t): 0\le x\le t<\infty\}$.
Разменяме реда на интегриране (позволено, защото $f_X\ge0$):
\[
  \int_0^{\infty}\Big(\int_0^t dx\Big) f_X(t)\,dt = \int_0^{\infty}
  t\,f_X(t)\,dt = \mathbb EX. \qedbg
\]
\end{proofbox}

\section*{Тема 10}

\begin{proofbox}[Решение на Задача 10.1]
За $X\sim Exp(\lambda)$: $\mathbb P(X>x)=e^{-\lambda x}$ за $x\ge0$.
Понеже $s+t\ge s$, $\{X>s+t\}\subseteq\{X>s\}$, следователно
\[
  \mathbb P(X>s+t\mid X>s) = \frac{\mathbb P(X>s+t)}{\mathbb P(X>s)}
  = \frac{e^{-\lambda(s+t)}}{e^{-\lambda s}} = e^{-\lambda t} =
  \mathbb P(X>t). \qedbg
\]
\end{proofbox}

\begin{proofbox}[Решение на Задача 10.2]
$g(x)=\tfrac{x-\mu}\sigma$ е строго растяща линейна функция с~обратна
$g^{-1}(z)=\sigma z+\mu$, $(g^{-1}(z))'=\sigma$. По теоремата за
смяна на променливите:
\[
  f_Z(z) = f_X(\sigma z+\mu)\cdot\sigma = \frac1{\sigma\sqrt{2\pi}}
  e^{-\frac{(\sigma z+\mu-\mu)^2}{2\sigma^2}}\cdot\sigma = \frac1{
  \sqrt{2\pi}} e^{-z^2/2},
\]
което е точно плътността на $N(0,1)$. \qedbg
\end{proofbox}

\begin{proofbox}[Решение на Задача 10.3]
$\mathbb P(X\le m)=1-e^{-\lambda m}=\tfrac12 \Rightarrow e^{-\lambda
m}=\tfrac12 \Rightarrow -\lambda m=-\ln2 \Rightarrow m=\dfrac{\ln2}
\lambda$.
\end{proofbox}

\section*{Тема 11}

\begin{proofbox}[Решение на Задача 11.1]
Търсената вероятност е лицето на областта $\{(x,y)\in(0,1)^2 : x+y<
1\}$ спрямо равномерната плътност (която е константа $1$ върху
единичния квадрат). Тази област е триъгълник с~върхове $(0,0),(1,0),
(0,1)$ с~лице $\tfrac12$, следователно $\mathbb P(X+Y<1)=\tfrac12$.
\end{proofbox}

\begin{proofbox}[Решение на Задача 11.2]
$Exp(\lambda)$ е частен случай на Гама разпределението: плътността
$\lambda e^{-\lambda x}$ съвпада с~плътността на $\Gamma(1,\lambda)$
при $\alpha=1$. Следователно $X_i\sim\Gamma(1,\lambda)$ независими,
и~по теоремата за адитивност на Гама разпределението (Тема~11):
\[
  \sum_{i=1}^n X_i \sim \Gamma(\underbrace{1+1+\dots+1}_{n\text{ пъти}
  },\ \lambda) = \Gamma(n,\lambda). \qedbg
\]
\end{proofbox}

\begin{proofbox}[Решение на Задача 11.3]
Пределни плътности: $f_X(x)=\int_0^1(x+y)\,dy = x+\tfrac12$ за
$x\in(0,1)$ (по симетрия, $f_Y(y)=y+\tfrac12$). Тогава
\[
  \mathbb EX = \int_0^1 x\Big(x+\frac12\Big)dx = \frac13+\frac14 =
  \frac7{12} = \mathbb EY \text{ (по симетрия)},
\]
\[
  \mathbb EXY = \int_0^1\!\!\int_0^1 xy(x+y)\,dy\,dx = \int_0^1\Big(
  \frac{x^2}2+\frac x3\Big)dx = \frac16+\frac16=\frac13.
\]
Следователно
\[
  \operatorname{cov}(X,Y) = \mathbb EXY-\mathbb EX\,\mathbb EY =
  \frac13 - \Big(\frac7{12}\Big)^2 = \frac{48}{144}-\frac{49}{144} =
  -\frac1{144}.
\]
\end{proofbox}

\section*{Тема 12}

\begin{proofbox}[Решение на Задача 12.1]
Фиксираме $\varepsilon>0$ и~полагаме $\delta=\min\big(1,\
\tfrac\varepsilon{1+2|c|}\big)>0$. Ако $|X_n-c|<\delta$, то $|X_n+c|
\le|X_n-c|+2|c| < 1+2|c|$, следователно
\[
  |X_n^2-c^2| = |X_n-c|\cdot|X_n+c| < \delta(1+2|c|) \le
  \varepsilon.
\]
Следователно $\{|X_n^2-c^2|>\varepsilon\}\subseteq\{|X_n-c|\ge
\delta\}$, откъдето
\[
  \mathbb P(|X_n^2-c^2|>\varepsilon) \le \mathbb P(|X_n-c|\ge\delta)
  \xrightarrow[n\to\infty]{}0,
\]
понеже $X_n\xrightarrow{\mathbb P}c$ и~$\delta>0$ е фиксирано.
Следователно $X_n^2\xrightarrow{\mathbb P}c^2$. \qedbg
\end{proofbox}

\begin{proofbox}[Решение на Задача 12.2]
По неравенството на Чебишов, приложено към извадковата средна (както
в~доказателството на СЗГЧ): $\mathbb P(|\overline X_n-\mu|>\varepsilon)
\le \dfrac{\sigma^2}{n\varepsilon^2}$. С $\sigma^2=4$, $\varepsilon=
0.5$, искаме
\[
  \frac{4}{n\cdot0.25} \le 0.05 \quad\Longleftrightarrow\quad
  \frac{16}n\le0.05 \quad\Longleftrightarrow\quad n\ge320.
\]
Следователно $n=320$ е достатъчно.
\end{proofbox}

\begin{proofbox}[Решение на Задача 12.3]
За Бернулиева величина, $DX_i=p(1-p)\le\tfrac14$ (максимумът на
$p(1-p)$ е при $p=\tfrac12$). По неравенството на Чебишов, приложено
към $\hat p_n=\overline X_n$:
\[
  \mathbb P(|\hat p_n-p|>0.02) \le \frac{DX_1}{n\cdot(0.02)^2} \le
  \frac{1/4}{10000\cdot0.0004} = \frac{1/4}{4} = \frac1{16} =
  0.0625.
\]
\end{proofbox}

\section*{Тема 13}

\begin{proofbox}[Решение на Задача 13.1]
$S_n\sim Bin(100,0.5)$, $np=50$, $\sqrt{np(1-p)}=\sqrt{25}=5$. По
Моавър--Лаплас, $\tfrac{S_n-50}5\xrightarrow d N(0,1)$. Следователно
\[
  \mathbb P(45<S_n<55) = \mathbb P\Big(-1<\frac{S_n-50}5<1\Big)
  \approx \Phi(1)-\Phi(-1) = 2\Phi(1)-1 \approx 2(0.8413)-1 =
  0.6826.
\]
\end{proofbox}

\begin{proofbox}[Решение на Задача 13.2]
$M_X(t)=\lambda(\lambda-t)^{-1}$. Първа производна: $M_X'(t)=
\lambda(\lambda-t)^{-2}$, при $t=0$: $\mathbb EX=M_X'(0)=\lambda/
\lambda^2=1/\lambda$. Втора производна: $M_X''(t)=2\lambda(\lambda-t)
^{-3}$, при $t=0$: $\mathbb EX^2=M_X''(0)=2\lambda/\lambda^3=2/
\lambda^2$. Следователно
\[
  DX = \mathbb EX^2-(\mathbb EX)^2 = \frac2{\lambda^2}-\frac1{
  \lambda^2} = \frac1{\lambda^2},
\]
което съвпада с~познатите $\mathbb EX=1/\lambda$, $DX=1/\lambda^2$
за експоненциалното разпределение.
\end{proofbox}

\begin{proofbox}[Решение на Задача 13.3]
За $X_1\sim Be(0.5)$: $\mu=0.5$, $\sigma^2=0.25$, $\sigma=0.5$, и~$|X_1-\mu|=0.5$ винаги (тъй като $X_1\in\{0,1\}$), следователно
$\mathbb E|X_1-\mu|^3 = 0.5^3=0.125$. С $N=100$:
\[
  \frac{0.4748\cdot0.125}{0.5^3\cdot\sqrt{100}} = \frac{0.4748\cdot
  0.125}{0.125\cdot10} = \frac{0.4748}{10} \approx 0.0475.
\]
\end{proofbox}

\section*{Тема 14}

\begin{proofbox}[Решение на Задача 14.1]
\[
  L(\bar X,\lambda) = \prod_{j=1}^n \frac{\lambda^{X_j}e^{-\lambda}}
  {X_j!} = \frac{\lambda^{\sum_jX_j}e^{-n\lambda}}{\prod_jX_j!},
  \qquad \ln L = \Big(\sum_jX_j\Big)\ln\lambda - n\lambda - \sum_j
  \ln(X_j!).
\]
\[
  \frac{\partial\ln L}{\partial\lambda} = \frac{\sum_jX_j}\lambda -
  n = 0 \quad\Longrightarrow\quad \hat\lambda = \frac1n\sum_{j=1}^n
  X_j = \overline X_n.
\]
\end{proofbox}

\begin{proofbox}[Решение на Задача 14.2]
$\mu^{(1)}(\lambda)=\mathbb EX=\dfrac1\lambda$. Изравняваме с~извадковия момент: $\dfrac1\lambda=\overline X_n \Rightarrow
\hat\lambda = \dfrac1{\overline X_n}$.
\end{proofbox}

\begin{proofbox}[Решение на Задача 14.3]
По неравенството на Чебишов, приложено към $\hat\theta_n$ (с~$\mathbb
E\hat\theta_n=\theta$ по неизместеност):
\[
  \mathbb P(|\hat\theta_n-\theta|>\varepsilon) = \mathbb P\big(|
  \hat\theta_n-\mathbb E\hat\theta_n|>\varepsilon\big) \le \frac{
  D\hat\theta_n}{\varepsilon^2} \xrightarrow[n\to\infty]{} 0,
\]
защото $D\hat\theta_n\to0$ по условие, а~$\varepsilon>0$ е
фиксирано. Понеже това важи за всяко $\varepsilon>0$, $\hat\theta_n
\xrightarrow{\mathbb P}\theta$, т.е.\ $\hat\theta_n$ е състоятелна.
\qedbg
\end{proofbox}

\section*{Тема 15}

\begin{proofbox}[Решение на Задача 15.1]
По ЦГТ, приложена към i.i.d.\ извадката $(X_i)$ с~$\mathbb EX_1=p$,
$DX_1=p(1-p)$:
\[
  T = \frac{\hat p_n-p}{\sqrt{p(1-p)/n}} \xrightarrow[n\to\infty]{d}
  N(0,1).
\]
Третирайки $T$ приблизително като централна статистика и~повтаряйки
конструкцията от Тема~15 със симетрични квантили $\pm q$, $q=
\Phi^{-1}(\tfrac{1+\gamma}2)$:
\[
  \gamma \approx \mathbb P\big(-q\le T\le q\big) = \mathbb P\Big(
  \hat p_n - q\sqrt{\tfrac{p(1-p)}n} \le p \le \hat p_n+q\sqrt{
  \tfrac{p(1-p)}n}\Big).
\]
Тъй като $p$ (в~корена) е неизвестно, на практика заместваме с~$\hat p_n$ (състоятелна оценка, оправдано защото $\hat p_n\to p$ по
вероятност), получавайки приблизителния доверителен интервал
\[
  \Big(\hat p_n - q\sqrt{\tfrac{\hat p_n(1-\hat p_n)}n},\ \ \hat p_n
  + q\sqrt{\tfrac{\hat p_n(1-\hat p_n)}n}\Big).
\]
\end{proofbox}

\begin{proofbox}[Решение на Задача 15.2]
По формулата за доверителен интервал при известно $\sigma^2$ (Тема
15), с~$\hat\mu=50$, $\sigma=2$, $n=16$, $q\approx1.96$:
\[
  \Big(50-1.96\cdot\frac2{\sqrt{16}},\ 50+1.96\cdot\frac2{\sqrt{16}}
  \Big) = \big(50-0.98,\ 50+0.98\big) = (49.02,\ 50.98).
\]
\end{proofbox}

\begin{proofbox}[Решение на Задача 15.3]
При $H_0$, тестовата статистика е
\[
  Z = \frac{\hat\mu-50}{\sigma/\sqrt n} = \frac{52-50}{4/\sqrt{25}} =
  \frac2{0.8} = 2.5.
\]
Критичната стойност е $q_{0.975}\approx1.96$. Понеже $|Z|=2.5>1.96$,
наблюдението попада в~критичната област, следователно \textbf{
отхвърляме} $H_0$ при ниво на значимост $\alpha=0.05$.
\end{proofbox}

\section*{Тема 16}

\begin{proofbox}[Решение на Задача 16.1]
$\bar x = \tfrac{1+2+3+4}4=2.5$, $\bar y=\tfrac{2+3+5+6}4=4$.
Отклонения от $\bar x$: $-1.5,-0.5,0.5,1.5$; съответните $y_k$:
$2,3,5,6$.
\[
  \hat\beta_1 = \frac{\sum y_k(x_k-\bar x)}{\sum(x_k-\bar x)^2} =
  \frac{2(-1.5)+3(-0.5)+5(0.5)+6(1.5)}{(-1.5)^2+(-0.5)^2+0.5^2+1.5^2}
  = \frac{-3-1.5+2.5+9}{2.25+0.25+0.25+2.25} = \frac7{5} = 1.4.
\]
\[
  \hat\beta_0 = \bar y - \bar x\,\hat\beta_1 = 4 - 2.5\cdot1.4 = 4-
  3.5 = 0.5.
\]
\end{proofbox}

\begin{proofbox}[Решение на Задача 16.2]
Оценките $(\hat\beta_0,\hat\beta_1)$ минимизират $\sum_{k=1}^n(Y_k-
b_0-b_1X_k)^2$ по $(b_0,b_1)$. Условието за минимум по $b_0$
(частна производна, приравнена на нула) е точно
\[
  \frac{\partial}{\partial b_0}\sum_{k=1}^n(Y_k-b_0-b_1X_k)^2
  \bigg|_{b_0=\hat\beta_0,\,b_1=\hat\beta_1} = -2\sum_{k=1}^n\big(
  Y_k-\hat\beta_0-\hat\beta_1X_k\big) = 0,
\]
откъдето директно $\sum_{k=1}^n(Y_k-\hat\beta_0-\hat\beta_1X_k)=0$.
\qedbg
\end{proofbox}

\begin{proofbox}[Решение на Задача 16.3]
По конструкцията от Тема~16, тестът отхвърля $H_0:\beta_1=b$ точно
когато
\[
  |\hat\beta_1-b| \ge q_{1-\alpha/2}\cdot\frac{\sigma}{\sqrt A},
  \qquad A=\sum_k(X_k-\bar X)^2.
\]
От друга страна, $(1-\alpha)$-доверителният интервал за $\beta_1$
(конструиран по same пивот $\tfrac{\hat\beta_1-\beta_1}{\sigma/\sqrt
A}\sim N(0,1)$, със същите симетрични квантили $\pm q_{1-\alpha/2}$)
е точно
\[
  \Big(\hat\beta_1-q_{1-\alpha/2}\frac\sigma{\sqrt A},\ \ \hat\beta_1
  +q_{1-\alpha/2}\frac\sigma{\sqrt A}\Big).
\]
Стойността $b$ \textbf{не} принадлежи на този интервал точно когато
$|\hat\beta_1-b|\ge q_{1-\alpha/2}\sigma/\sqrt A$ - същото
неравенство като условието за отхвърляне на $H_0$. Следователно
двете твърдения са еквивалентни. \qedbg
\end{proofbox}
